% Pursuing Stacks (cotes 134-2 to 134-8): archive cote 134-7, site page 6 (PDF page 7), typescript page 448.
% typescript. Draft compared: ps6.tex:232-295.
% AI-transcribed by Claude Opus 5.5 (claude-opus-5-5, adaptive thinking, effort high), Message Batches API, batch ps3 (msgbatch_01XdFmuVVBAcgc7EFxdtCqsj), 2026-11-12; not yet checked by eye.
\scanpage{134-7/6}
\mspage{448}

\deleted{\illegible{1}}
sheaves \xout{\unsure{of sets}} (of sets) just described.  Possibly both approaches are
of interest,  and will yield non equivalent notions of schematization.
On the other hand,  although definitely X is not representable by a
usual \xout{\unsure{affine}} scheme over k unless the k-modules $M_i$ are projective of
finite type (in which case X will be an affine scheme,  and even isomor-%
phic,  at least locally over Spec(k),  to standard affine space $E_k^d$ for
suitable  d),  it is felt that X,  as far as cohomology properties go,
should be very close to being an affine scheme,  and that presumably
it's cohomology groups $H^i$  with coefficients in ``quasi-coherent'' \xout{\illegible{1}}
sheaves such as W(M),  should vanish for $i>0$;  consequently,  presumably
the torsors used for the inductive construction of X are trivial,  which
means that X is in fact isomorphic to the product of all $W(M_i)$'s .  In
the case when we disregard the successive fibration structure,  this
means that the objects of U(k) are just sheaves of sets which are isomor-%
phic to some W(M)  (where morally M is the direct sum of the modules $M_i$
which\xout{\illegible{1}} have been used in our inductive definition).  This gives then a
rather trivial description of the objects of U(k) (and all $U_m(k)$'s are
already equal to U(k)\hand{,} for $m \over{>}{\geq} 1$ ! ),  it should be remembered however \xout{\unsure{that}}
that maps in U(k)  from a W(M) to a W(M')  are a lot more general than
just k-linear maps  $M \to M'$  (they may be viewed as ``polynomial maps
from M to M'\,'').

The category U(k) will be endowed with the \uline{sections} \uline{functor}
\[ \text{(h)} \qquad X \mapsto X(k) \ : \ U(k) \longrightarrow \text{(Sets)} \quad . \]
Now let A be any test category,  for instance A = \hand{$\Delta$},  and consider
the functor
\[ \hand{(h')} \qquad \underline{\mathrm{Hom}}(A^{o},U(k)) \longrightarrow \underline{\mathrm{Hom}}(A^{o},\text{(Sets)}) \]
induced by h).  The second member modelizes homotopy types,  which there-%
fore allows us to defined homotopy invariants for objects in the first
member,  and hence to define the property of n-connectedness,  and (if
n=0)  of 0-connectedness with abelian fundamental group.  As a matter of
fact,  we would like to define a subcategory $\hand{$\mathcal{M}$}_n(k)$   of the first mem-%
ber,  so that it should become clear that for an object $X_{\hand{$*$}}$  in it, it's
homotopy groups are endowed with structures of k-modules,  so that for
n=0  the abelian restriction on $\hand{$\pi$}_1$ is superfluous,  because automatic.
To be specific,  we better restrict maybe to A= \hand{$\Delta$},  more familiar to
us,  so that $X_{\hand{$*$}}$ may be viewed as a ``complex''  ($n \mapsto X_n$)  with components
in U(k).  The kind of restriction I am thinking of for defining the

\margin{6}
